%
%Academic CV LaTeX Template
% Author: Dubasi Pavan Kumar
% LinkedIn: https://www.linkedin.com/in/im-pavankumar/
% License: MIT
%
% For errors, suggestions, or improvements, please contact:
% Email: pavankumard.pg19.ma@nitp.ac.in
%



\documentclass[a4paper,11pt]{article}

% Package imports
\usepackage{latexsym}
\usepackage{xcolor}
\usepackage{float}
\usepackage{ragged2e}
\usepackage[empty]{fullpage}
\usepackage{wrapfig}
\usepackage{lipsum}
\usepackage{tabularx}
\usepackage{titlesec}
\usepackage{geometry}
\usepackage{marvosym}
\usepackage{verbatim}
\usepackage{enumitem}
\usepackage{fancyhdr}
\usepackage{multicol}
\usepackage{graphicx}
\usepackage{cfr-lm}
\usepackage[T1]{fontenc}
\usepackage{fontawesome5}
% Vietnamese language support
\usepackage[utf8]{inputenc}
\usepackage[T5]{fontenc}
\usepackage[vietnamese]{babel}

% Color definitions
\definecolor{darkblue}{RGB}{0,0,139}

% Page layout
\setlength{\multicolsep}{0pt} 
\pagestyle{fancy}
\fancyhf{} % clear all header and footer fields
\fancyfoot{}
\renewcommand{\headrulewidth}{0pt}
\renewcommand{\footrulewidth}{0pt}
\geometry{left=1.4cm, top=0.8cm, right=1.2cm, bottom=1cm}
\setlength{\footskip}{5pt} % Addressing fancyhdr warning

% Hyperlink setup (moved after fancyhdr to address warning)
\usepackage[hidelinks]{hyperref}
\hypersetup{
    colorlinks=true,
    linkcolor=darkblue,
    filecolor=darkblue,
    urlcolor=darkblue,
}

% ORCID support
\usepackage{orcidlink}

% Custom box settings
\usepackage[most]{tcolorbox}
\tcbset{
    frame code={},
    center title,
    left=0pt,
    right=0pt,
    top=0pt,
    bottom=0pt,
    colback=gray!20,
    colframe=white,
    width=\dimexpr\textwidth\relax,
    enlarge left by=-2mm,
    boxsep=4pt,
    arc=0pt,outer arc=0pt,
}

% URL style
\urlstyle{same}

% Text alignment
\raggedright
\setlength{\tabcolsep}{0in}

% Section formatting
\titleformat{\section}{
  \vspace{-4pt}\scshape\raggedright\large
}{}{0em}{}[\color{black}\titlerule \vspace{-7pt}]

% Custom commands
\newcommand{\resumeItem}[2]{
  \item{
    \textbf{#1}{\hspace{0.5mm}#2 \vspace{-0.5mm}}
  }
}

\newcommand{\resumePOR}[3]{
\vspace{0.5mm}\item
    \begin{tabular*}{0.97\textwidth}[t]{l@{\extracolsep{\fill}}r}
        \textbf{#1}\hspace{0.3mm}#2 & \textit{\small{#3}} 
    \end{tabular*}
    \vspace{-2mm}
}

\newcommand{\resumeSubheading}[4]{
\vspace{0.5mm}\item
    \begin{tabular*}{0.98\textwidth}[t]{l@{\extracolsep{\fill}}r}
        \textbf{#1} & \textit{\footnotesize{#4}} \\
        \textit{\footnotesize{#3}} &  \footnotesize{#2}\\
    \end{tabular*}
    \vspace{-2.4mm}
}

\newcommand{\resumeProject}[4]{
\vspace{0.5mm}\item
    \begin{tabular*}{0.98\textwidth}[t]{@{\extracolsep{\fill}}>{\raggedright\arraybackslash}p{0.76\textwidth}r}
        \textbf{#1} & \textit{\footnotesize{#3}} \\
        \footnotesize{\textit{#2}} & \footnotesize{#4}
    \end{tabular*}
    \vspace{-2.4mm}
}

\newcommand{\resumeSubItem}[2]{\resumeItem{#1}{#2}\vspace{-4pt}}

\renewcommand{\labelitemi}{$\vcenter{\hbox{\tiny$\bullet$}}$}
\renewcommand{\labelitemii}{$\vcenter{\hbox{\tiny$\circ$}}$}

\newcommand{\resumeSubHeadingListStart}{\begin{itemize}[leftmargin=*,labelsep=1mm]}
\newcommand{\resumeHeadingSkillStart}{\begin{itemize}[leftmargin=*,itemsep=1.7mm, rightmargin=2ex]}
\newcommand{\resumeItemListStart}{\begin{itemize}[leftmargin=*,labelsep=1mm,itemsep=0.5mm]}

\newcommand{\resumeSubHeadingListEnd}{\end{itemize}\vspace{2mm}}
\newcommand{\resumeHeadingSkillEnd}{\end{itemize}\vspace{-2mm}}
\newcommand{\resumeItemListEnd}{\end{itemize}\vspace{-2mm}}
\newcommand{\cvsection}[1]{%
\vspace{2mm}
\begin{tcolorbox}
    \textbf{\large #1}
\end{tcolorbox}
    \vspace{-4mm}
}

\newcolumntype{L}{>{\raggedright\arraybackslash}X}%
\newcolumntype{R}{>{\raggedleft\arraybackslash}X}%
\newcolumntype{C}{>{\centering\arraybackslash}X}%

% Commands for icon sizing and positioning
\newcommand{\socialicon}[1]{\raisebox{-0.05em}{\resizebox{!}{1em}{#1}}}
\newcommand{\ieeeicon}[1]{\raisebox{-0.3em}{\resizebox{!}{1.3em}{#1}}}

% Font options
\newcommand{\headerfonti}{\fontfamily{phv}\selectfont} % Helvetica-like (similar to Arial/Calibri)
\newcommand{\headerfontii}{\fontfamily{ptm}\selectfont} % Times-like (similar to Times New Roman)
\newcommand{\headerfontiii}{\fontfamily{ppl}\selectfont} % Palatino (elegant serif)
\newcommand{\headerfontiv}{\fontfamily{pbk}\selectfont} % Bookman (readable serif)
\newcommand{\headerfontv}{\fontfamily{pag}\selectfont} % Avant Garde-like (similar to Trebuchet MS)
\newcommand{\headerfontvi}{\fontfamily{cmss}\selectfont} % Computer Modern Sans Serif
\newcommand{\headerfontvii}{\fontfamily{qhv}\selectfont} % Quasi-Helvetica (another Arial/Calibri alternative)
\newcommand{\headerfontviii}{\fontfamily{qpl}\selectfont} % Quasi-Palatino (another elegant serif option)
\newcommand{\headerfontix}{\fontfamily{qtm}\selectfont} % Quasi-Times (another Times New Roman alternative)
\newcommand{\headerfontx}{\fontfamily{bch}\selectfont} % Charter (clean serif font)

\begin{document}
\headerfontiii

% Header
\begin{center}
    {\Huge\textbf{VƯƠNG VIỆT HÙNG}}
\end{center}
\vspace{-5mm}

\begin{center}
    \small{
    +84-583-885-391 | \href{mailto:vuongviethung156@gmail.com}{vuongviethung156@gmail.com}
    }
\end{center}
\vspace{-5mm}

\begin{center}
    \small{
    \socialicon{\faLinkedin} \href{https://www.linkedin.com/in/v%C6%B0%C6%A1ng-vi%E1%BB%87t-h%C3%B9ng-843199212/}{vương-việt-hùng} | 
    \socialicon{\faGithub} \href{https://github.com/vuong-viet-hung}{vuong-viet-hung} | \orcidlink{0009-0006-8041-7608} \href{https://orcid.org/my-orcid?orcid=0009-0006-8041-7608}{0009-0006-8041-7608}
    }
\end{center}
\vspace{-5mm}
\begin{center}
    \small{Hanoi, Vietnam - 11025, Vietnam}
\end{center}

\vspace{-4mm}

\section{\textbf{Education}}
\vspace{-0.4mm}
\resumeSubHeadingListStart

\resumeSubheading
{Soongsil University}{Seoul, South Korea}
{Master's Degree in Electronic Engineering}{March 2024 - March 2026}
\resumeItemListStart
\item GPA: 4.24/4.50
\resumeItemListEnd

\resumeSubheading
{Hanoi University of Science and Technology}{Hanoi, Vietnam}
{Bachelor's Degree in Control Engineering and Automation}{August 2019 - September 2023}
\resumeItemListStart
\item GPA: 3.34/4.00
\resumeItemListEnd

\resumeSubHeadingListEnd
\vspace{-6mm}

\section{\textbf{Experience}}

\vspace{-0.4mm}

\resumeSubHeadingListStart

\resumeSubheading
  {VinDynamics}{Hanoi, Vietnam}
  {SLAM \& Navigation Engineer}{March 2026 -- current}
\vspace{3mm}
\small{Developed and deployed SLAM and height mapping systems for humanoid robotic platforms, with a focus on sensor integration, system evaluation, and optimization.}

\resumeSubheading
  {ANDA Lab, Soongsil University}{Seoul, South Korea}
  {Research Assistant}{March 2024 -- March 2026}
\vspace{3mm}
\small{Conducted research in LiDAR panoptic segmentation, scene completion, domain adaptation, and continual learning; contributed to peer-reviewed publications, researcher mentoring, and technical seminars.}

\resumeSubheading
  {Rikai Mind}{Hanoi, Vietnam}
  {AI Engineer Intern}{August 2022 -- January 2023}
\vspace{3mm}
\small{Developed and deployed a computer vision system for object volume estimation from monocular images and camera--object distance.}

\resumeSubheading
  {Mandevices Laboratory, Hanoi University of Science and Technology}{Hanoi, Vietnam}
  {Research Assistant}{August 2020 -- September 2023}
\vspace{3mm}
\small{Conducted research in deep-learning-based bearing fault diagnosis and contributed to data collection, publications, and researcher mentoring.}

\resumeSubHeadingListEnd
\vspace{-6mm}


\section{\textbf{Publications}}

\resumeSubHeadingListStart

  \resumeProject
    {Domain Adaptive LiDAR Panoptic Segmentation via Pseudo-Supervised Instance Learning}
    {First author -- IEEE Transactions on Intelligent Transportation Systems, Early Access}
    {May 2026}

  {\small DOI: \href{https://doi.org/10.1109/TITS.2026.3695516}{10.1109/TITS.2026.3695516}}

  \resumeItemListStart
    \item Proposed PILLA, a domain adaptation framework that generates instance-level pseudo-labels via unsupervised clustering and applies a noise-aware loss for robust pseudo-supervision.
    \item Achieved 35.97\% PQ on nuScenes in cross-dataset adaptation from SemanticKITTI, outperforming AdaptLPS by 13.14 percentage points.
  \resumeItemListEnd


  \resumeProject
    {Instance Embedding as Queries for DETR-based LiDAR Panoptic Segmentation}
    {Co-first author -- IEEE Transactions on Intelligent Vehicles, Vol. 10, Issue 10, pp. 4629 - 4641}
    {October 2025}

  {\small DOI: \href{https://doi.org/10.1109/TIV.2024.3488035}{10.1109/TIV.2024.3488035}}

  \resumeItemListStart
    \item Proposed a DETR-based architecture that uses instance embeddings as object queries for LiDAR panoptic segmentation.
    \item Achieved state-of-the-art performance with 63.40\% PQ on SemanticKITTI benchmark and 77.08\% PQ on nuScenes benchmark, outperforming previous methods by 1.20 and 1.94 percentage points, respectively.
  \resumeItemListEnd


  \resumeProject
    {FuPaSCo: Long-range and Local Context Fusion for 3D Panoptic Scene Completion}
    {Co-author -- Image and Vision Computing, Vol. 163, pp. 105776}
    {November 2025}

  {\small DOI: \href{https://doi.org/10.1016/j.imavis.2025.105776}{10.1016/j.imavis.2025.105776}}

  \resumeItemListStart
    \item Developed FuPaSCo, a dual-branch CNN--Transformer framework that combines local and long-range features through feature fusion and alignment for LiDAR panoptic scene completion.
    \item Achieved 26.10\% $PQ^\dagger$ on SemanticKITTI benchmark, outperforming the previous best method PasCo by 2.49 percentage points.
  \resumeItemListEnd


  \resumeProject
    {Generalized Simulation-Based Domain Adaptation Approach for Intelligent Bearing Fault Diagnosis}
    {Co-author -- Arabian Journal for Science and Engineering, Vol. 49, Issue 12, pp. 16941--16957}
    {July 2024}

  {\small DOI: \href{https://doi.org/10.1007/s13369-024-09282-1}{10.1007/s13369-024-09282-1}}

  \resumeItemListStart
    \item Developed a simulation-based domain adaptation approach that generates synthetic vibration signals from bearing geometry to address data scarcity in intelligent bearing fault diagnosis.
    \item Achieved 74--98\% fault classification accuracy across different bearing types using a CNN trained on synthesized vibration signals.
  \resumeItemListEnd


  \resumeProject
    {ConvMamba: A Data-Efficient Neural Network for Bearing Fault Diagnosis}
    {Co-author -- 7th International Conference on Green Technology and Sustainable Development (GTSD), pp. 165-172. IEEE. 2024-07-25, Ho Chi Minh City, Vietnam}
    {July 2024}

  {\small DOI: \href{https://doi.org/10.1109/GTSD62346.2024.10674880}{10.1109/GTSD62346.2024.10674880}}

  \resumeItemListStart
    \item Developed ConvMamba, a data-efficient architecture combining convolutional networks and Vision Mamba for bearing fault diagnosis.
    \item Improved classification accuracy by 27 percentage points over the predecessor Hybrid-CNN-MLP model on the CWRU bearing fault dataset.
  \resumeItemListEnd

\resumeSubHeadingListEnd


\section{\textbf{Manuscripts Under Review}}

\resumeSubHeadingListStart

\resumeProject
  {SM3-PSC: Sparse Multi-Modal Mamba for 3D Panoptic Scene Completion with Multi-Scale Feature Refinement}
  {Co-author -- Manuscript submitted to IEEE Transactions on Image Processing}
  {2026}

\resumeItemListStart
  \item Developed SM3-PSC, a sparse multi-modal framework that fuses LiDAR and camera features and employs Mamba-based multi-scale feature refinement for 3D panoptic scene completion.
  \item Achieved 29.20\% $PQ^\dagger$ on SemanticKITTI, outperforming the previous best method by 3.10 percentage points through improved completion of both object instances and surrounding regions.
\resumeItemListEnd


\resumeProject
  {LW-PSC: Lightweight Design for LiDAR Panoptic Scene Completion}
  {Co-author -- Manuscript submitted to Applied Soft Computing}
  {2026}

\resumeItemListStart
  \item Developed LW-PSC, a lightweight BEV-based framework that combines Cartesian and Polar representations with center-focused feature refinement for real-time LiDAR panoptic scene completion.
  \item Achieved 13.59\% PQ at 11.3 FPS on SemanticKITTI, outperforming PaSCo by 1.47 percentage points in PQ while reducing the model size from 40.65M to 29.16M parameters.
\resumeItemListEnd

\resumeSubHeadingListEnd


\section{\textbf{Projects}}

\resumeSubHeadingListStart

  \resumeProject
    {Visual SLAM Profiling and Latency Optimization}
    {VinDynamics}
    {August 2026}

  \resumeItemListStart
    \item Profiled the OKVIS2 visual SLAM pipeline, identifying landmark matching and local-window optimization as major bottlenecks, and evaluated low-latency front-end pose publishing with only 1.16 cm translational APE RMSE relative to back-end poses.

    \item Redesigned frame scheduling to prioritize the latest available frame and drop stale frames, reducing average front-end latency from 124.00 ms to 47.00 ms and back-end latency from 164.30 ms to 87.64 ms.
  \resumeItemListEnd


  \resumeProject
    {Multi-Sensor Local Height Mapping for Humanoid Robot}
    {VinDynamics}
    {July 2026 -- August 2026}

  \resumeItemListStart
    \item Extended a LiDAR-based GPU height mapper to integrate Livox LiDAR and Intel RealSense depth data, implementing sensor synchronization, rate limiting, pose lookup, and robot self-filtering for deployment on an in-house humanoid robot.

    \item Implemented and benchmarked pose extrapolation to compensate for SLAM latency, reducing translational APE RMSE from 2.148 cm to 0.445 cm (79.3\%) for a 30 Hz depth camera using linear extrapolation.
  \resumeItemListEnd


  \resumeProject
    {Scene Flow-Guided Online Continual LiDAR Panoptic Tracking}
    {Graduation thesis, Soongsil University}
    {September 2025 -- February 2026}

  \resumeItemListStart

    \item Proposed SF-LPT, the first online continual learning framework for LiDAR panoptic tracking, using self-supervised scene flow for continuous adaptation and temporal instance association across sequential scans.
    \item Achieved 16.45\% LSTQ on SemanticKITTI in cross-dataset continual adaptation from nuScenes, outperforming CoDEPS by 5.87 percentage points.

  \resumeItemListEnd


  \resumeProject
    {Stereo Matching and 3D Reconstruction}
    {Graduation thesis, Hanoi University of Science and Technology}
    {March 2023 -- August 2023}

  \resumeItemListStart
    \item Implemented and compared block matching and deep learning approaches for dense stereo matching on SceneFlow (EPE: 77.32 px, D1: 79.04\%).
    \item Reconstructed 3D scenes from estimated disparity maps via triangulation and deployed the pipeline on a Jetson Nano-based wheeled robot.
  \resumeItemListEnd


  \resumeProject
    {Object Volume Estimation from Images}
    {Rikai Mind}
    {August 2022 -- December 2022}

  \resumeItemListStart
    \item Developed an object volume estimation model from monocular images and camera--object distance, achieving an $R^2$ score of 84.78\%.
    \item Extracted visual embeddings with CNNs, localized objects using keypoint regression, and fused visual features with camera--object distance for volume regression; deployed the model using FastAPI.
  \resumeItemListEnd

\resumeSubHeadingListEnd


\section{\textbf{Skills}}
\vspace{-0.4mm}
 \resumeHeadingSkillStart
  \resumeSubItem{Programming languages: }
    {Python, C/C++, MATLAB}
  \resumeSubItem{Frameworks: }
    {PyTorch, TensorFlow, OpenCV, ROS2, Kalibr (for sensor calibration)}
  \resumeSubItem{Hardware platforms: }
    {Jetson Orin Nano, Jetson Nano, Raspberry Pi}
  \resumeSubItem{Sensors: }
    {Intel RealSense camera, ZED X camera, Livox LiDAR}
  \resumeSubItem{Development and deployment tools: }
    {Git, GitHub, GitLab, Docker, FastAPI, Milvus}
  \resumeSubItem{Others: }{LaTex}
 \resumeHeadingSkillEnd

\section{\textbf{Languages}}
\vspace{-0.4mm}
 \resumeHeadingSkillStart
  \resumeSubItem{English: }
    {IELTS 7.5}
  \resumeSubItem{Vietnamese: }
    {Native speaker}
 \resumeHeadingSkillEnd

\end{document}